\chapter{综合结果与综合后仿真}\label{ch:netlist}

使用 eLinx 对 \code{AFDX_top} 进行综合，并导出 VQM 网表。综合后网表通过双向通信测试（NET-01）和五项功能联合测试（NET-02）。

\section{综合配置与资源}

综合采用 eLinx~3.0.7 工具，顶层模块为 \code{AFDX_top}，目标器件选用 EQ6HL130（CSG484\_H 封装），优化策略设为 Speed，主时钟 \code{clk} 约束为 125~MHz（周期 8~ns）。综合流程无错误完成并导出结构化 VQM 网表，整机资源消耗汇总见表~\ref{tab:test-synth-res}。

\par\medskip\noindent\begin{minipage}{\linewidth}\renewcommand{\tablearraystretch}{1.16}\tablestyle\centering
\captionof{table}{总体综合资源}\label{tab:test-synth-res}
\begin{tabularx}{\linewidth}{LRRR}
\toprule
\thead{资源} & \thead{用量} & \thead{器件总量} & \thead{占比} \\
\midrule
LUT & \Lut & 136160 & \LutPct\% \\
FF & \FF & 136160 & \FFPct\% \\
CLB & \Clb & 13616 & \ClbPct\% \\
存储块（BRAM） & \Mem & 1200 & \MemPct\% \\
DSP & 0 & 48 & 0\% \\
PLL & 0 & 8 & 0\% \\
用户引脚 & \Pins & 339 & \PinPct\% \\
\bottomrule
\end{tabularx}
\end{minipage}\par\medskip

表~\ref{tab:test-synth-mod} 依据综合层次报告列出顶层下除轻量级计数器 \code{afdx_timebase} 外的 7 个主要功能模块，以及收发核心内部的关键控制子模块（各父模块资源均包含其下属子模块及片上存储实例）。发送与接收核心合计使用 14491 个逻辑单元和 2315136 位片上存储，集中承载 VL 配置表、调度与冗余状态、描述符队列及报文数据缓存。

\begingroup\renewcommand{\tablearraystretch}{1.16}
\begin{longtable}{@{}>{\raggedright\arraybackslash}p{5.4cm}>{\raggedleft\arraybackslash}p{2.6cm}>{\raggedleft\arraybackslash}p{2.6cm}>{\raggedleft\arraybackslash}p{\dimexpr\linewidth-10.6cm-6\tabcolsep\relax}@{}}
\caption{主要模块资源}\label{tab:test-synth-mod}\\
\longtablehead{4}{\thead{模块} & \thead{逻辑单元} & \thead{FF} & \thead{存储位数}}
\code{afdx_tx_core} & 5677 & 2975 & 1106816 \\
\quad \code{afdx_tx_engine} & 1790 & 1029 & 0 \\
\quad \code{afdx_tx_builder} & 2091 & 962 & 0 \\
\quad \code{afdx_mac_tx} & 597 & 155 & 0 \\
\code{afdx_rx_core} & 8814 & 4964 & 1208320 \\
\quad \code{afdx_rx_parse}（2 个） & 2466 & 1220 & 0 \\
\quad \code{afdx_rx_engine} & 1127 & 489 & 0 \\
\quad \code{afdx_rx_deliver} & 1975 & 481 & 0 \\
\quad \code{afdx_mac_rx}（2 个） & 371 & 222 & 0 \\
\code{afdx_rtc} & 4573 & 2190 & 0 \\
\code{afdx_snmp} & 4463 & 1261 & 4096 \\
\code{afdx_615a} & 7171 & 1358 & 102400 \\*
\code{afdx_cfg} & 1954 & 1076 & 0 \\*
\code{afdx_svc} & 120 & 7 & 0 \\
\end{longtable}
\endgroup

表中“逻辑单元”来自层次报告的 Entity 统计口径（含仅使用寄存器的逻辑单元，故与表~\ref{tab:test-synth-res} 的纯 LUT 统计口径不同）；缩进子模块已计入所属核心，标注“2 个”的行代表双网通道两个实例的合计值。“存储位数”统计综合后实际映射的有效数据位宽：其中 \code{afdx_rx_core} 为 1208320 位，比设计文档表 9-2 按 RTL 数组位宽声明统计的 1228800 位少 20480 位，原因是综合器自动裁除了接收配置表 \code{rx_cfg}（$4096\times32$~bit）每项中未使用的最高 5 位（$4096\times5=20480$~bit）。


\subsection{MAC 单口资源对照}

在相同目标器件 EQ6HL130 下，建立两个独立的单口工程分别对本设计自研的流式千兆 MAC（\code{afdx_mac_tx}+\code{afdx_mac_rx}）与 eLinx IP 核生成器导出的 TripleSpeedEthernet IP（含用户侧 FIFO 的标准千兆配置）进行综合打包对比，结果见表~\ref{tab:test-mac-res}（原始打包报告见附录~\ref{app:evidence}，整机双网集成资源见表~\ref{tab:test-synth-res}）。自研 MAC 省去了通用 IP 中的跨时钟域 FIFO、管理接口及冗余速率适配逻辑，单口 LUT 与寄存器开销分别仅为通用 IP 的 39.7\% 和 20.3\%，且不占用片上存储块。

\par\medskip\noindent\begin{minipage}{\linewidth}\renewcommand{\tablearraystretch}{1.16}\tablestyle\centering
\captionof{table}{MAC 单口资源对照}\label{tab:test-mac-res}
\begin{tabularx}{\linewidth}{LRR}
\toprule
\thead{资源} & \thead{自研 MAC} & \thead{TripleSpeedEthernet IP 配置} \\
\midrule
LUT & 556 & 1400 \\
FF & 198 & 975 \\
存储块（BRAM） & 0 & 24 \\
\bottomrule
\end{tabularx}
\end{minipage}\par\medskip

\section{综合时序}

在主时钟 \code{clk} 周期 8~ns（125~MHz）约束下，eLinx 综合阶段的静态时序分析摘要见表~\ref{tab:test-synth-time}。综合阶段采用逻辑单元延迟与预估时钟树模型，数据路径的互连布线延迟按 0~ns 计算。

\par\medskip\noindent\begin{minipage}{\linewidth}\renewcommand{\tablearraystretch}{1.16}\tablestyle\centering
\captionof{table}{8~ns（125~MHz）时钟约束下的综合时序摘要}\label{tab:test-synth-time}
\begin{tabularx}{\linewidth}{@{}>{\raggedright\arraybackslash}p{3.0cm}>{\raggedleft\arraybackslash}p{2.8cm}>{\raggedleft\arraybackslash}p{2.8cm}L@{}}
\toprule
\thead{检查项} & \thead{综合报告值} & \thead{违例端点/\newline 检查端点} & \thead{说明} \\
\midrule
最高频率估计 & \Fmax~MHz & — & 高于 125~MHz 目标频率 \\
Setup 最差裕量 & +\Setup~ns & 0/44216 & 综合模型下满足 8~ns 周期约束 \\
Hold 最差裕量 & $-$\Hold~ns & 538/44216 & 零布线延迟模型下的检查结果 \\
\bottomrule
\end{tabularx}
\end{minipage}\par\medskip

全设计 44216 个同步端点均满足 8~ns 建立时间（Setup）约束，工具评估的最高工作频率为 \Fmax~MHz。建立时间裕量最小的路径位于 \code{afdx_615a} 模块内部，起点为 \code{name_pos[1]}、终点为 \code{packet_type[1]}，经 15 级组合逻辑，总延迟 5.403~ns，保留 +\Setup~ns 正裕量。

在零互连延迟的综合前模型下，保持时间（Hold）检查报告 538 个跨存储块端点的表面负裕量，最差值为 $-$\Hold~ns（位于发送缓存写数据寄存器 \code{payload_wd} 至 BRAM 写端口路径）。该现象源于综合模型中 BRAM 时钟树分支比普通触发器晚 0.754~ns，而数据路径仅计入 0.696~ns 的触发器输出单元延迟、互连布线延迟按 0~ns 计，从而产生 68~ps 的名义缺口；进入布局布线阶段计入数百皮秒的实际互连延迟与绕线补偿后即可自然闭合。

\Needspace{16\baselineskip}
\section{网表功能结果}

综合后仿真直接复用 RTL 测试平台与原厂单元功能仿真库。由于 eLinx 导出的结构化 VQM 网表省略了原语未用输入端口的默认电平绑定，仿真前通过脚本 \code{tools/adapt_gui_vqm.py} 按原厂原语声明显式补齐悬空输入的默认电平与进位链使能参数，同时完整保留网表原有的全部信号连接、配置参数及 LUT 掩码值（详见附录~\ref{sec:repro-netlist}）。

NET-01 执行双端双向通信测试平台自检；NET-02 在完成自检与独立参考模型核对的基础上，进一步将网表仿真的外部端口记录与 RTL 仿真记录进行含精确时钟周期编号的逐字节比对。两项测试均一次性通过~\ref{sec:criteria}~节规定的全部检查，且 ModelSim 仿真全过程无 Error、Fatal 或 Warning。

\begin{items}{综合后仿真测试项}{tab:net-result}
NET-01 & 由界面综合导出的网表；激励同 SYS-09（\code{tb_elinx}） & 自检通过；仿真诊断无 Error、Fatal、Warning & \ElinxChecks/0\newline 平台自检 & 通过 \\
NET-02 & 由界面综合导出的网表；激励同 SVC-08（\code{tb_net_services}） & 自检和独立核对通过；7 份端口记录与 RTL 逐字节相同；仿真诊断无 Error、Fatal、Warning & \SvcChecks/0\newline 65/0 & 通过 \\
\end{items}

NET-02 的 7 份端口记录与对应的 RTL 记录逐字节一致，包括五项功能联合测试总记录、各端各网的线路记录、接收记录和卸载结果。加载并卸载的 16 字节文件与原文件一致，卸载内容的 41 项检查全部通过。
